\section{国家赔偿}
\begin{enumerate}
  \item 国家进行赔偿，包括行政赔偿和司法赔偿（民事、刑事）。钱从国库出。
  \item \zd{混合归责}原则。
        \begin{enumerate}
          \item \zd{过错归责}。事实行为。同民法。
          \item \zd{违法归责}。只看事中不管事后。
          \item \zd{结果归责}。不看事中\zd{看事后}。例外。拘留违法规则，\zd{超期拘留}（>= 37 天）转结果。
        \end{enumerate}
  \item 构成要件。主体、行为、结果、因果关系。
  \item 行政赔偿。
        \begin{kk}{赔什么、谁来赔、赔给谁、怎么赔}
          \begin{enumerate}
            \item 赔什么。侵害人身权。\al{唆使牢头狱霸放纵他人殴打虐待}。
            \item 谁来赔（叫赔偿义务人）。确定方式类似诉讼被告。
                  \begin{enumerate}
                    \item 复议加重侵害\zd{按份}，各赔各。
                    \item \zd{派出所永远不赔}。
                    \item 做\zd{基础决定}的行政机关。
                    \item 共同侵权（无意思联络的数人侵权）共同赔。原因力\zd{竞合/结合 \yjt 连带/按份}。
                    \item 第三人侵权。\al{安保义务人}。
                  \end{enumerate}
            \item 赔给谁。同原告，\zd{支付了医疗费和丧葬费等合理费用}的人\zd{可以}要求赔偿。
            \item 怎么赔。
                  \begin{enumerate}
                    \item \zd{一并赔}。在复议/诉讼程序中直接申请。
                    \item \zd{单独赔}。2 年内。先行处理（谁赔找谁）/其他方式确认违法 \yjt 提起行政赔偿诉讼。
                  \end{enumerate}
          \end{enumerate}
          \tcblower
          \begin{enumerate}
            \item 行政诉讼、行政复议、\zd{行政赔偿诉讼}三者平起平坐。
            \item \zd{举证责任}。谁主张谁举证。例外。\zd{限制人身自由 + 发生损害结果}。几乎必考。
            \item \zd{不可抗力一般不赔}。自己作用力范围内。
          \end{enumerate}
        \end{kk}
  \item 司法赔偿。
        \begin{enumerate}
          \item 人身自由。
                \begin{kk}{没罪关了就要赔，有罪关了也白关；免罪关了分前后，赔后不赔前}
                  \begin{enumerate}
                    \item 没罪。真的无罪、疑罪从无。
                    \item 关。实际羁押，\zd{减刑部分、保外就医、管制、缓期执行、取保候审}不赔。
                    \item 免罪。显著微，过时效，特救告诉和死掉。
                    \item 前后。以\zd{生效判决生效}为界限。
                  \end{enumerate}
                \end{kk}
          \item 原则上谁做谁赔，例外时后置原则（\zd{谁最后做有罪决定谁赔}）。
                \begin{enumerate}
                  \item 公检赔。3 步走。自己 \yjt 爹 \yjt 爹的同级法院。
                  \item 法院赔。两步走。自己 \yjt 爹。
                \end{enumerate}
          \item 程序。\zd{没有确认违法程序了}。
                \begin{enumerate}
                  \item 可以口头提，时效 \zd{2} 年。
                  \item 非诉程序。\zd{书面不开庭}。\zd{叫申请赔偿}，选项出现\zd{赔偿诉讼}一律不选！！！！！！！！。
                  \item 赔偿委员会；\zd{3} 人单数。
                  \item 救济。\zd{申诉而非申请复议}，就一个重新处理程序，类似再审。
                \end{enumerate}
          \item 金额计算。\zd{决定赔偿时上年度职工日平均工资}。1000 起算（坚持要 1 块 \yjt 解释 \yjt 仍坚持 \yjt 准！）。
                \begin{enumerate}
                  \item 人身自由按日支付；误工费 <= 5 倍，丧葬费 20 倍。
                  \item 半残没有生活费。\zd{1-4 全残，5-10 半残}。
                        \begin{kk}{精神赔偿。关了打了才要赔，伤害轻了还不赔。}
                          \begin{enumerate}
                            \item 有告知受害人有权请求赔偿的义务。
                            \item 一般严重。超过 6 个月、轻伤、半残。总额 50\% 以下。
                            \item 特别严重。10 年以上、重伤、全残、死亡。总额 50\% 以上。
                          \end{enumerate}
                        \end{kk}
                \end{enumerate}
          \item 财产类。
                \begin{enumerate}
                  \item 侵权行为时市场价格。
                  \item 变卖要给差额，\zd{拍卖拍多少赔多少}。
                  \item 只赔\zd{直接损失/必然损失}。\al{房租、折旧费、贷款利息、机动车营运、应得补助 √；预期利润 ×}。
                \end{enumerate}
        \end{enumerate}
\end{enumerate}
